\documentclass[11pt,a4paper]{article}
\usepackage{amsmath,amssymb}
\usepackage[margin=2.5cm]{geometry}

\title{Mamba-3 Paper Equation~(2):\\Derivation of the Matrix Form of MIMO SSM}
\author{}
\date{}

\begin{document}
\maketitle

We derive the output matrix $\mathbf{Y} = (\mathbf{L} \odot \mathbf{CB}^\top) \mathbf{X}$ in Mamba-3, starting from the fundamental equations of the continuous-time system and proceeding step by step with explicit dimensions.

\section{Definitions: Dimensions and Parameters}

\subsection{Basic Dimensions}

\begin{itemize}
  \item $L$: Sequence Length
  \item $D$: Number of Input/Output Channels (Model Dimension)
  \item $N$: State Dimension (per channel)
\end{itemize}

\subsection{Variable Dimension Definitions}

\begin{itemize}
  \item $\mathbf{X} \in \mathbb{R}^{L \times D}$: Input sequence
  \item $\mathbf{Y} \in \mathbb{R}^{L \times D}$: Output sequence
  \item $h(t) \in \mathbb{C}^{D \times N}$: State matrix. Each column $h_d(t) \in \mathbb{C}^{N \times 1}$ is the state vector of channel~$d$.
  \item $\mathbf{A}(t)$: Collection of diagonal components.
    \begin{itemize}
      \item Storage dimension: $\mathbb{C}^{D \times N}$. Each element $A_{d,n}(t)$ is the complex coefficient for the $n$-th state of channel~$d$.
      \item Operational dimension: For each channel~$d$, it acts as an $N \times N$ diagonal matrix $\mathbf{A}_d(t) = \mathrm{diag}(A_{d,1}(t), \dots, A_{d,N}(t))$.
    \end{itemize}
  \item $\mathbf{B}(t) \in \mathbb{R}^{L \times N}$, $\mathbf{C}(t) \in \mathbb{R}^{L \times N}$: Projection vectors shared across all channels.
\end{itemize}

\section{Proof: Expansion from Continuous Time to Matrix Form}

\subsection*{Step 1: Derivation of the Continuous-Time Solution (Integrating Factor Method)}

We begin with the differential equation for the state $h_d(t) \in \mathbb{C}^{N \times 1}$ (column vector) of each channel~$d$:
\begin{equation}
  \frac{dh_d(t)}{dt} = \mathbf{A}_d(t)\, h_d(t) + \mathbf{b}(t)\, x_d(t)
\end{equation}
where $\mathbf{A}_d(t) \in \mathbb{C}^{N \times N}$ is a diagonal matrix, $\mathbf{b}(t) \in \mathbb{R}^{N \times 1}$, and $x_d(t) \in \mathbb{R}$.

To solve this first-order linear non-homogeneous differential equation, we multiply both sides by the integrating factor $\exp\!\bigl(-\int_{0}^{t} \mathbf{A}_d(\tau)\,d\tau\bigr)$:
\begin{equation}
  \exp\!\Bigl(-\!\int_{0}^{t}\!\mathbf{A}_d(\tau)\,d\tau\Bigr)
  \biggl[\frac{dh_d(t)}{dt} - \mathbf{A}_d(t)\,h_d(t)\biggr]
  = \exp\!\Bigl(-\!\int_{0}^{t}\!\mathbf{A}_d(\tau)\,d\tau\Bigr)\,\mathbf{b}(t)\,x_d(t)
\end{equation}

We apply the reverse of the product rule $\frac{d}{dt}(fg) = f'g + fg'$ to the left-hand side:
\begin{equation}
  \frac{d}{dt}\biggl[\exp\!\Bigl(-\!\int_{0}^{t}\!\mathbf{A}_d(\tau)\,d\tau\Bigr)\,h_d(t)\biggr]
  = \exp\!\Bigl(-\!\int_{0}^{t}\!\mathbf{A}_d(\tau)\,d\tau\Bigr)\,\mathbf{b}(t)\,x_d(t)
\end{equation}

Setting the initial state $h_d(0)=0$ and integrating both sides from $0$ to $t$:
\begin{equation}
  \exp\!\Bigl(-\!\int_{0}^{t}\!\mathbf{A}_d(\tau)\,d\tau\Bigr)\,h_d(t)
  - \underbrace{\exp(0)\,h_d(0)}_{=\,0}
  = \int_{0}^{t}\exp\!\Bigl(-\!\int_{0}^{s}\!\mathbf{A}_d(\tau)\,d\tau\Bigr)\,\mathbf{b}(s)\,x_d(s)\,ds
\end{equation}

Multiplying both sides from the left by $\exp\!\bigl(\int_{0}^{t}\mathbf{A}_d(\tau)\,d\tau\bigr)$:
\begin{equation}
  h_d(t) = \exp\!\Bigl(\int_{0}^{t}\!\mathbf{A}_d(\tau)\,d\tau\Bigr)
  \int_{0}^{t}\exp\!\Bigl(-\!\int_{0}^{s}\!\mathbf{A}_d(\tau)\,d\tau\Bigr)\,\mathbf{b}(s)\,x_d(s)\,ds
\end{equation}

Since $\mathbf{A}$ is a diagonal matrix, we can bring the exponential term inside the integral:
\begin{equation}
  h_d(t) = \int_{0}^{t}\exp\!\Bigl(\int_{0}^{t}\!\mathbf{A}_d(\tau)\,d\tau - \int_{0}^{s}\!\mathbf{A}_d(\tau)\,d\tau\Bigr)\,\mathbf{b}(s)\,x_d(s)\,ds
\end{equation}

By the linearity of integration $\int_{0}^{t} - \int_{0}^{s} = \int_{s}^{t}$, we obtain the continuous-time general solution:
\begin{equation}\label{eq:continuous}
  h_d(t) = \int_{0}^{t}
  \underbrace{\exp\!\Bigl(\int_{s}^{t}\!\mathbf{A}_d(\tau)\,d\tau\Bigr)}_{\displaystyle\Phi_d(t,s)}
  \,\mathbf{b}(s)\,x_d(s)\,ds
  \qquad\bigl[h_d(t)\in\mathbb{C}^{N \times 1}\bigr]
\end{equation}

\subsection*{Step 2: Discretization and Derivation of the Cumulative Product}

\paragraph{2.1\quad Discrete Update Rule and Definition of $\bar{\mathbf{b}}_k$.}

We derive the state $h_{d,k}$ at time $t = k\Delta$ with time step~$\Delta$ by splitting the integration range $[0,\,k\Delta]$ into $[0,\,(k{-}1)\Delta]$ and $[(k{-}1)\Delta,\,k\Delta]$:
\begin{equation}
  h_{d,k} = \int_{0}^{(k-1)\Delta}\!\Phi_d(k\Delta,s)\,\mathbf{b}(s)\,x_d(s)\,ds
           + \int_{(k-1)\Delta}^{k\Delta}\!\Phi_d(k\Delta,s)\,\mathbf{b}(s)\,x_d(s)\,ds
\end{equation}

We apply the semigroup property $\Phi_d(t,s)=\Phi_d(t,u)\,\Phi_d(u,s)$ to the first term:
\begin{equation}
  h_{d,k} = \Phi_d\bigl(k\Delta,(k{-}1)\Delta\bigr)
  \underbrace{\int_{0}^{(k-1)\Delta}\!\Phi_d\bigl((k{-}1)\Delta,s\bigr)\,\mathbf{b}(s)\,x_d(s)\,ds}_{\displaystyle h_{d,k-1}}
  + \int_{(k-1)\Delta}^{k\Delta}\!\Phi_d(k\Delta,s)\,\mathbf{b}(s)\,x_d(s)\,ds
\end{equation}

We define the discretization parameters as follows:
\begin{itemize}
  \item Transition matrix:
    $\displaystyle\mathbf{A}_{d,k} = \Phi_d\bigl(k\Delta,(k{-}1)\Delta\bigr) = \exp\!\Bigl(\int_{(k-1)\Delta}^{k\Delta}\!\mathbf{A}_d(\tau)\,d\tau\Bigr)$
  \item Discretized input projection:
    $\displaystyle\bar{\mathbf{b}}_k = \int_{(k-1)\Delta}^{k\Delta}\!\Phi_d(k\Delta,s)\,\mathbf{b}(s)\,ds$
    \quad(assuming $x_d(s)$ is constant within the interval)
\end{itemize}

Using these, we obtain the discrete update rule:
\begin{equation}\label{eq:scalar-update}
  \underbrace{h_{d,k}}_{N\times 1}
  = \underbrace{\mathbf{A}_{d,k}}_{N\times N}\;\underbrace{h_{d,k-1}}_{N\times 1}
  + \underbrace{\bar{\mathbf{b}}_k}_{N\times 1}\;\underbrace{x_{d,k}}_{1\times 1}
\end{equation}

\paragraph{2.2\quad Extension to Matrix Form (All-Channel Integration).}

We now derive the matrix form (paper equation~11) from the per-channel update~\eqref{eq:scalar-update}.

\medskip\noindent\textbf{Collecting all channels into a matrix.}\quad
Define the state matrix $\mathbf{H}_k \in \mathbb{C}^{D\times N}$ whose $d$-th row is $h_{d,k}^\top$:
\begin{equation}
  \mathbf{H}_k
  = \begin{pmatrix} h_{1,k}^\top \\ h_{2,k}^\top \\ \vdots \\ h_{D,k}^\top \end{pmatrix}
  \in \mathbb{C}^{D\times N}
\end{equation}

\medskip\noindent\textbf{Rewriting the transition term.}\quad
Since $\mathbf{A}_{d,k} = \mathrm{diag}(A_{d,k,1},\dots,A_{d,k,N})$ is diagonal, the matrix--vector product reduces to an elementwise product:
\begin{equation}
  \bigl(\mathbf{A}_{d,k}\,h_{d,k-1}\bigr)_n = A_{d,k,n}\cdot h_{d,k-1,n}
\end{equation}
Transposing to row form, the $d$-th row of the transition term is:
\begin{equation}
  \bigl(\mathbf{A}_{d,k}\,h_{d,k-1}\bigr)^\top
  = \bigl(A_{d,k,1}\,h_{d,k-1,1},\;\dots,\;A_{d,k,N}\,h_{d,k-1,N}\bigr)
\end{equation}
Define $\mathcal{A}_k \in \mathbb{C}^{D\times N}$ with $(\mathcal{A}_k)_{d,n} = A_{d,k,n}$, i.e.\ the matrix that collects all diagonal entries. Stacking all $D$ rows:
\begin{equation}
  \begin{pmatrix}
    (\mathbf{A}_{1,k}\,h_{1,k-1})^\top \\
    \vdots \\
    (\mathbf{A}_{D,k}\,h_{D,k-1})^\top
  \end{pmatrix}
  = \mathcal{A}_k \odot \mathbf{H}_{k-1}
\end{equation}
where $\odot$ denotes the Hadamard (elementwise) product.

\medskip\noindent\textbf{Rewriting the input term.}\quad
The input contribution for channel~$d$ is the column vector $\bar{\mathbf{b}}_k\,x_{d,k}$. Transposing to row form gives $x_{d,k}\,\bar{\mathbf{b}}_k^\top \in \mathbb{R}^{1\times N}$. Stacking all $D$ rows with $\mathbf{x}_k = (x_{1,k},\dots,x_{D,k})^\top \in \mathbb{R}^{D\times 1}$:
\begin{equation}
  \begin{pmatrix}
    x_{1,k}\,\bar{\mathbf{b}}_k^\top \\
    \vdots \\
    x_{D,k}\,\bar{\mathbf{b}}_k^\top
  \end{pmatrix}
  = \mathbf{x}_k\,\bar{\mathbf{b}}_k^\top
  \in \mathbb{R}^{D\times N}
\end{equation}
This is an outer product: the $(d,n)$-entry is $x_{d,k}\,\bar{b}_{k,n}$.

\medskip\noindent\textbf{Combining both terms.}\quad
Summing the transition and input terms yields the matrix update rule:
\begin{equation}\label{eq:matrix-update}
  \boxed{\;\mathbf{H}_k = \mathcal{A}_k \odot \mathbf{H}_{k-1} + \mathbf{x}_k\,\bar{\mathbf{b}}_k^\top\;}
\end{equation}

\paragraph{2.3\quad Expansion by Recursive Substitution.}

By substituting recursively starting from $h_0 = 0$, the state at any time~$t$ is:
\begin{equation}
  h_{d,t} = \sum_{s=0}^{t}\Biggl(\prod_{j=s+1}^{t}\mathbf{A}_{d,j}\Biggr)\bar{\mathbf{b}}_s\,x_{d,s}
\end{equation}

\subsection*{Step 3: Substitution into the Output Equation and Kernel Extraction}

Substituting into the output $y_{d,t} = \mathbf{c}_t \cdot h_{d,t}$ and expanding componentwise using the properties of diagonal matrices:
\begin{equation}
  y_{d,t} = \sum_{s=0}^{t}\Biggl[\sum_{n=1}^{N}(c_{t,n}\cdot\bar{b}_{s,n})\prod_{j=s+1}^{t}A_{d,j,n}\Biggr]x_{d,s}
\end{equation}

\subsection*{Step 4: Consolidation into Matrix Form (Paper Equation~2)}

Define:
\begin{itemize}
  \item Transition kernel $\mathbf{L}\in\mathbb{C}^{L\times L\times D}$:\quad
    $(t,s,d)$-element $\displaystyle L_{ts,d}=\prod_{j=s+1}^{t}A_{d,j}$
  \item Projection kernel $\mathbf{C}\mathbf{B}^\top\in\mathbb{R}^{L\times L}$:\quad
    $(t,s)$-element $(\mathbf{C}\mathbf{B}^\top)_{ts}=\mathbf{c}_t^\top\bar{\mathbf{b}}_s$
\end{itemize}

Then:
\begin{equation}\label{eq:main}
  \boxed{\;\mathbf{Y} = \bigl(\mathbf{L}\odot\mathbf{C}\mathbf{B}^\top\bigr)\,\mathbf{X}\;}
\end{equation}

\hfill$\blacksquare$

\end{document}
